% ANEXOS

% --- INÍCIO DA SEÇÃO DE APÊNDICES ---
\appendix
\chapter{Apêndices: Informações Adicionais}
\label{chap:anexos}

\section{Estrutura de Arquivos do Projeto}

O projeto está organizado da seguinte forma:

\begin{itemize}
    \item \texttt{codes/}: Contém todos os scripts Python implementados
    \begin{itemize}
        \item \texttt{2A-prepare\_corel\_dataset.py}: Preparação do dataset
        \item \texttt{2B-train-lora-corel.py}: Treinamento LoRA
        \item \texttt{2C-generate-lora-corel.py}: Geração com LoRA
        \item \texttt{3A-train-vae-corel.py}: Treinamento VAE
        \item \texttt{3B-train-diffusion-from-scratch-corel.py}: Treinamento difusão
        \item \texttt{3C-generate-samples-corel.py}: Geração com difusão
        \item \texttt{4A-train-dgae-corel.py}: Treinamento DGAE
        \item \texttt{4B-extract-features-corel.py}: Extração features DGAE
        \item \texttt{5A-simclr\_corel.py}: Treinamento SimCLR
        \item \texttt{5D-cnn-jepa\_corel.py}: Treinamento CNN-JEPA
        \item \texttt{5E-compare-clustering.py}: Comparação de técnicas
    \end{itemize}
    \item \texttt{README\_2.md}, \texttt{README\_3.md}, \texttt{README\_4.md}, \texttt{README\_5.md}: Documentação de cada tarefa
    \item \texttt{training\_data/corel/}: Dataset preparado
    \item \texttt{features/}: Features extraídas para clustering
    \item \texttt{clustering\_results/}: Resultados de comparação
\end{itemize}

\section{Parâmetros de Treinamento Principais}

\subsection{LoRA (Item 2)}

\begin{itemize}
    \item Base model: \texttt{runwayml/stable-diffusion-v1-5}
    \item LoRA rank: 4
    \item LoRA alpha: 4
    \item Learning rate: 1e-4
    \item Batch size: 1
    \item Gradient accumulation: 4
    \item Epochs: 200 (configurável)
    \item Resolution: 512x512
\end{itemize}

\subsection{VAE (Item 3)}

\begin{itemize}
    \item Image size: 128x128
    \item Latent dimension: 128
    \item Hidden dimensions: [64, 128, 256, 512]
    \item Learning rate: 1e-4
    \item Batch size: 16
    \item KL weight: 1.0
    \item KL warmup: 150 epochs
    \item Perceptual weight: 0.03
    \item Epochs: 300
\end{itemize}

\subsection{Difusão Latente (Item 3)}

\begin{itemize}
    \item Learning rate: 1e-4
    \item Batch size: 16
    \item Gradient accumulation: 1
    \item Image size: 128x128 (após VAE)
    \item EMA decay: 0.995
    \item Schedule: linear
    \item Mixed precision: FP16
    \item Epochs: 500
\end{itemize}

\subsection{DGAE (Item 4)}

\begin{itemize}
    \item Image size: 256x256
    \item Latent dimension: 128
    \item Hidden dimensions: [64, 128, 256, 512]
    \item Learning rate: 1e-4
    \item Batch size: 8
    \item Diffusion guidance weight: 0.1
    \item Perceptual weight: 0.03
    \item Epochs: 300
\end{itemize}

\subsection{SimCLR (Item 5)}

\begin{itemize}
    \item Image size: 224x224
    \item Latent dimension: 128
    \item Learning rate: 3e-3
    \item Temperature: 0.5
    \item Batch size: 32
    \item Epochs: 100
\end{itemize}

\subsection{CNN-JEPA (Item 5)}

\begin{itemize}
    \item Image size: 224x224
    \item Feature dimension: 512
    \item Hidden dimension: 512
    \item Learning rate: 1e-3
    \item Batch size: 32
    \item Epochs: 100
\end{itemize}

\section{Comandos de Execução}

\subsection{Item 2: LoRA}

\begin{verbatim}
# Preparar dataset
python 2A-prepare_corel_dataset.py \
    --data-dir data/corel --base-dir .

# Treinar LoRA
python 2B-train-lora-corel.py \
    --data-dir training_data/corel/corel_all \
    --output-dir corel_models \
    --epochs 200

# Gerar imagens
python 2C-generate-lora-corel.py \
    --lora-dir corel_models \
    --num-images 16
\end{verbatim}

\subsection{Item 3: Difusão do Zero}

\begin{verbatim}
# Treinar VAE
python 3A-train-vae-corel.py \
    --data-dir training_data/corel/corel_all \
    --output-dir vae_models \
    --epochs 300

# Treinar difusão
python 3B-train-diffusion-from-scratch-corel.py \
    --vae-checkpoint vae_models/best_model.pt \
    --image-dir training_data/corel/corel_all \
    --output-dir diffusion_models

# Gerar amostras
python 3C-generate-samples-corel.py \
    --diffusion-checkpoint diffusion_models/best_model.pt \
    --vae-checkpoint vae_models/best_model.pt \
    --num-images 16
\end{verbatim}

\subsection{Item 4: DGAE}

\begin{verbatim}
# Treinar DGAE
python 4A-train-dgae-corel.py \
    --data-dir training_data/corel/corel_all \
    --lora-dir corel_models \
    --output-dir dgae_models

# Extrair features
python 4B-extract-features-corel.py \
    --model-checkpoint dgae_models/best_model.pt \
    --data-dir training_data/corel/corel_all \
    --output features/dgae_features.npy \
    --extract-labels
\end{verbatim}

\subsection{Item 5: Clustering}

\begin{verbatim}
# Treinar SimCLR
python 5A-simclr_corel.py \
    --data-dir training_data/corel/corel_all \
    --output-dir simclr_models \
    --extract-features features/simclr_features.npy

# Treinar CNN-JEPA
python 5D-cnn-jepa_corel.py \
    --data-dir training_data/corel/corel_all \
    --output-dir cnn_jepa_models \
    --extract-features features/cnn_jepa_features.npy

# Comparar técnicas
python 5E-compare-clustering.py \
    --features-dir features \
    --output-dir clustering_results \
    --compare-both
\end{verbatim}

\section{Métricas de Clustering}

As métricas utilizadas para avaliação são definidas como:

\subsection{Adjusted Rand Index (ARI)}

O ARI mede a concordância entre os clusters obtidos e as classes verdadeiras, ajustado para o acaso:

\[
ARI = \frac{\sum_{ij} \binom{n_{ij}}{2} - [\sum_i \binom{a_i}{2} \sum_j \binom{b_j}{2}] / \binom{n}{2}}{\frac{1}{2}[\sum_i \binom{a_i}{2} + \sum_j \binom{b_j}{2}] - [\sum_i \binom{a_i}{2} \sum_j \binom{b_j}{2}] / \binom{n}{2}}
\]

onde $n_{ij}$ é o número de elementos na classe $i$ e cluster $j$, $a_i$ é o número de elementos na classe $i$, e $b_j$ é o número de elementos no cluster $j$.

\subsection{Normalized Mutual Information (NMI)}

O NMI mede a informação mútua normalizada entre clusters e classes:

\[
NMI = \frac{2 \cdot I(C;K)}{H(C) + H(K)}
\]

onde $I(C;K)$ é a informação mútua entre classes $C$ e clusters $K$, e $H(\cdot)$ é a entropia.

\subsection{Silhouette Score}

O Silhouette Score mede quão similar um objeto é ao seu próprio cluster comparado a outros clusters:

\[
s(i) = \frac{b(i) - a(i)}{\max(a(i), b(i))}
\]

onde $a(i)$ é a distância média aos outros pontos no mesmo cluster e $b(i)$ é a distância média aos pontos no cluster mais próximo.

\section{Arquiteturas de Modelos}

\subsection{Encoder SimCLR/CNN-JEPA}

\begin{itemize}
    \item Conv2d(3, 64, 7x7, stride=2) + BN + ReLU
    \item MaxPool2d(3x3, stride=2)
    \item Conv2d(64, 128, 3x3, stride=2) + BN + ReLU
    \item Conv2d(128, 256, 3x3, stride=2) + BN + ReLU
    \item Conv2d(256, 512, 3x3, stride=2) + BN + ReLU
    \item AdaptiveAvgPool2d(1)
    \item Projection head (SimCLR) ou Linear (CNN-JEPA)
\end{itemize}

\subsection{VAE/DGAE Encoder}

\begin{itemize}
    \item Múltiplas camadas convolucionais com blocos residuais
    \item Downsampling progressivo (stride=2)
    \item Dimensões ocultas: [64, 128, 256, 512]
    \item Saída: Espaço latente de dimensão configurável (128 para DGAE)
\end{itemize}

\subsection{VAE/DGAE Decoder}

\begin{itemize}
    \item Input linear do espaço latente
    \item Upsampling progressivo com PixelShuffle
    \item Blocos residuais em cada nível
    \item Saída: Imagem reconstruída (3 canais RGB)
\end{itemize}

\section{Considerações de Implementação}

\subsection{Otimizações GPU}

Todos os scripts implementam:

\begin{itemize}
    \item Verificação de disponibilidade CUDA
    \item \texttt{cudnn.benchmark = True} para convoluções mais rápidas
    \item \texttt{pin\_memory = True} no DataLoader
    \item \texttt{num\_workers > 0} para carregamento paralelo
    \item Mixed precision (FP16) quando aplicável
    \item Limpeza periódica de cache CUDA
\end{itemize}

\subsection{Compatibilidade de Checkpoints}

\begin{itemize}
    \item Configurações salvas como dicionários (não objetos) para evitar erros de pickle
    \item Suporte para carregar checkpoints antigos e novos
    \item Inferência automática de parâmetros se config não disponível
\end{itemize}

\subsection{Reprodutibilidade}

\begin{itemize}
    \item Seeds fixos para randomização (padrão: 42)
    \item Ordenação consistente de arquivos
    \item Determinismo quando possível (com trade-off de performance)
\end{itemize}
